\clearpage
\seccionInforme{2. Diseño de interfaces}
\phantomsection\addcontentsline{toc}{section}{2. Diseño de interfaces}

\vspace{0.2cm}
\phantomsection\addcontentsline{toc}{subsection}{2.1 Funcionalidades derivadas del árbol de objetivos}
\colabtitle{2.1 Funcionalidades derivadas del árbol de objetivos}

Cada medio fundamental del árbol de objetivos se descompuso en funcionalidades concretas. La Tabla 4 muestra esa trazabilidad, que es la base del diseño: ninguna pantalla se agregó sin un medio que la justifique.

\begin{table}[H]
\centering\small
\caption{Trazabilidad entre medios del árbol de objetivos, funcionalidades e interfaces}
\renewcommand{\arraystretch}{1.25}
\begin{tabular}{|>{\raggedright\arraybackslash}p{\dimexpr 0.240\linewidth-2\tabcolsep-2\arrayrulewidth\relax}|>{\raggedright\arraybackslash}p{\dimexpr 0.530\linewidth-2\tabcolsep-2\arrayrulewidth\relax}|>{\raggedright\arraybackslash}p{\dimexpr 0.230\linewidth-2\tabcolsep-2\arrayrulewidth\relax}|}
\hline
\rowcolor{colabpyblue} \color{white}\textbf{Medio} & \color{white}\textbf{Funcionalidad} & \color{white}\textbf{Interfaz} \\ \hline
1.1 Catálogo conectado & F-01 Iniciar sesión con correo institucional y elegir rol & I1 Inicio de sesión \\ \hline
\rowcolor{colabpybg} 1.1 Catálogo conectado & F-02 Ver resumen personal: ventana, créditos, vacantes y laboratorios matriculados & I2 Panel del estudiante \\ \hline
1.1 Catálogo conectado & F-03 Explorar grupos por curso con semáforo, buscador y filtros & I3 Catálogo \\ \hline
\rowcolor{colabpybg} 2.1 Contraste de horario & F-04 Revisar el detalle del grupo y el mapa de puestos; validar cruces antes de confirmar & I4 Detalle del grupo \\ \hline
2.1 Contraste de horario & F-05 Consultar el horario semanal con teoría y laboratorios & I6 Horario \\ \hline
\rowcolor{colabpybg} 3.1 Operación única & F-06 Confirmar la matrícula con reserva atómica del cupo & I4 Detalle del grupo \\ \hline
4.1 Código verificable & F-07 Mostrar y compartir el pase con QR y código de verificación & I5 Pase QR \\ \hline
\rowcolor{colabpybg} 4.1 Código verificable & F-08 Monitorear la ocupación de los grupos del docente & I7 Panel docente \\ \hline
4.1 Código verificable & F-09 Consultar la nómina, marcar asistencia y exportar CSV & I8 Nómina \\ \hline
\rowcolor{colabpybg} 4.1 Código verificable & F-10 Leer el QR con la cámara o digitar el código & I9 Escáner \\ \hline
5.1 Parametrización & F-11 Ajustar capacidad, estado y datos de cada laboratorio & I10 Laboratorios \\ \hline
\rowcolor{colabpybg} 5.1 Parametrización & F-12 Configurar la ventana de matrícula y las reglas; revisar grupos agotados & I11 Matrícula \\ \hline
Transversal & F-13 Elegir tema claro u oscuro, ajustar avisos y cerrar sesión & I12 Perfil \\ \hline
\end{tabular}
\end{table}

\vspace{0.2cm}
\phantomsection\addcontentsline{toc}{subsection}{2.2 Mapa de navegación}
\colabtitle{2.2 Mapa de navegación}

La estructura de navegación se definió antes de dibujar las pantallas (Figura 3). Una \textbf{pila nativa} raíz contiene el inicio de sesión, tres navegadores de \textbf{pestañas inferiores} (uno por rol) y tres pantallas que se apilan sobre las pestañas: el detalle del grupo, que se presenta como hoja modal, el pase QR y la nómina. Así cada rol ve solo sus opciones y el botón Atrás del sistema se comporta como el usuario espera.

\begin{figure}[H]
\centering
\safeimage{1.0\textwidth}{img/diagramas/navegacion.png}{Mapa de navegación de MatriLab: pila nativa raíz y pestañas por rol}
\caption{Mapa de navegación de MatriLab: pila nativa raíz y pestañas por rol}
\label{fig:navegacion}
\end{figure}

\vspace{0.2cm}
\phantomsection\addcontentsline{toc}{subsection}{2.3 Sistema visual}
\colabtitle{2.3 Sistema visual}

Para que las doce interfaces se vean consistentes se definieron \textit{tokens} de diseño en \texttt{src/\allowbreak{}theme/\allowbreak{}tokens.\allowbreak{}ts}. Se conservaron los colores institucionales usados desde el LAB 01 y se definió una variante oscura de cada color, ya que la aplicación sigue el tema del sistema operativo.

\begin{table}[H]
\centering\small
\caption{Tokens de diseño de la aplicación}
\renewcommand{\arraystretch}{1.25}
\begin{tabular}{|>{\raggedright\arraybackslash}p{\dimexpr 0.200\linewidth-2\tabcolsep-2\arrayrulewidth\relax}|>{\raggedright\arraybackslash}p{\dimexpr 0.270\linewidth-2\tabcolsep-2\arrayrulewidth\relax}|>{\raggedright\arraybackslash}p{\dimexpr 0.200\linewidth-2\tabcolsep-2\arrayrulewidth\relax}|>{\raggedright\arraybackslash}p{\dimexpr 0.330\linewidth-2\tabcolsep-2\arrayrulewidth\relax}|}
\hline
\rowcolor{colabpyblue} \color{white}\textbf{Token} & \color{white}\textbf{Claro} & \color{white}\textbf{Oscuro} & \color{white}\textbf{Uso} \\ \hline
primario & \#19375A & \#7FA7D9 & Botones principales, cabeceras de tarjetas, grupo seleccionado. \\ \hline
\rowcolor{colabpybg} marca & \#C3141E & \#FF6B6B & Acciones decisivas (Ingresar, Confirmar matrícula) y pestaña activa. \\ \hline
éxito / alerta / peligro & \#1E7B45 / \#A35F00 / \#B3261E & variantes claras & Semáforo de vacantes: verde, ámbar y rojo. \\ \hline
\rowcolor{colabpybg} fondo / superficie & \#F4F6F9 / \#FFFFFF & \#0F1720 / \#17212C & Fondo de pantalla y de tarjetas. \\ \hline
espacio & 4 · 8 · 12 · 16 · 24 · 32 & igual & Escala de 4 pt para márgenes, rellenos y \texttt{gap}. \\ \hline
\rowcolor{colabpybg} radio & 8 · 12 · 16 · 999 & igual & Esquinas de controles, tarjetas y píldoras. \\ \hline
tipo & 26 · 18 · 15 · 13 · 12 & igual & Título, subtítulo, cuerpo, etiqueta y pie (en dp). \\ \hline
\end{tabular}
\end{table}

\vspace{0.2cm}
\phantomsection\addcontentsline{toc}{subsection}{2.4 Mockups de baja fidelidad}
\colabtitle{2.4 Mockups de baja fidelidad}

Con las funcionalidades y la navegación definidas se dibujaron mockups de baja fidelidad para las doce interfaces (Figuras 4 y 5). Se usaron solo grises, sin textos reales, para discutir la \textbf{organización} de cada pantalla antes que su apariencia. Cada región lleva una etiqueta roja con el componente de React Native previsto y los recuadros discontinuos marcan contenedores de layout (filas, columnas o listas). Ese mismo boceto guió luego la implementación: en la sección 4.6 cada mockup aparece junto a la captura real de la pantalla.

\begin{figure}[H]
\centering
\safeimage{0.8\textwidth}{img/diagramas/mockups_estudiante.png}{Mockups de baja fidelidad de las interfaces del estudiante (I1 a I6)}
\caption{Mockups de baja fidelidad de las interfaces del estudiante (I1 a I6)}
\label{fig:mockups_e}
\end{figure}

\begin{figure}[H]
\centering
\safeimage{0.8\textwidth}{img/diagramas/mockups_docente_admin.png}{Mockups de baja fidelidad de las interfaces del docente y de administración (I7 a I12)}
\caption{Mockups de baja fidelidad de las interfaces del docente y de administración (I7 a I12)}
\label{fig:mockups_d}
\end{figure}

Las decisiones de organización más importantes que quedaron fijadas en los mockups fueron las siguientes:

\begin{itemize}[label=\textcolor{unsared}{\textbullet}, leftmargin=*]
    \item \textbf{Listas largas virtualizadas.} El catálogo, el panel docente, la nómina y los laboratorios se planificaron como \texttt{Section\allowbreak{}List} o \texttt{Flat\allowbreak{}List}, nunca como un \texttt{Scroll\allowbreak{}View} con muchos hijos.
    \item \textbf{Acción principal fija.} En el detalle del grupo y en la nómina el botón principal queda en un pie fuera del área desplazable, siempre visible.
    \item \textbf{Superposiciones con posición absoluta.} El sello del pase, la guía de encuadre del escáner y los bloques del horario se colocan con \texttt{position: 'absolute'} sobre su contenedor.
    \item \textbf{Cuadrículas que se reacomodan.} Los indicadores (KPI) y el mapa de puestos usan \texttt{flex\allowbreak{}Wrap}, de modo que el número de columnas depende del ancho disponible y no de valores fijos.
\end{itemize}

\vspace{0.2cm}
\phantomsection\addcontentsline{toc}{subsection}{2.5 Criterios de usabilidad y accesibilidad}
\colabtitle{2.5 Criterios de usabilidad y accesibilidad}

\begin{itemize}[label=\textcolor{unsared}{\textbullet}, leftmargin=*]
    \item Área táctil mínima de \textbf{48 dp} en botones y campos (\texttt{min\allowbreak{}Height: 48}), y \texttt{hit\allowbreak{}Slop} en iconos pequeños.
    \item Todos los controles tienen \texttt{accessibility\allowbreak{}Label} y \texttt{accessibility\allowbreak{}Role}; las filas de la nómina se anuncian como casillas (\texttt{accessibility\allowbreak{}Role=\allowbreak{}"checkbox"}) con su estado.
    \item Los mensajes de error y de validación usan \texttt{accessibility\allowbreak{}Live\allowbreak{}Region} para que el lector de pantalla los lea al aparecer.
    \item El semáforo nunca depende solo del color: siempre va acompañado del texto \textit{Disponible}, \textit{Últimos cupos} o \textit{Agotado}.
    \item Tema claro y oscuro con contraste suficiente; la escala de texto del sistema (\texttt{font\allowbreak{}Scale}) se respeta y se muestra en el perfil.
\end{itemize}

